% Chapter — Time systems (FR-2, D-6). Follows the mandatory FR-29
% six-section template: purpose; assumptions and domain of validity;
% derivation; discretization and implementation; validation evidence;
% references. The equation labels eq:time:j2000, eq:time:fvf,
% eq:time:utc2tai, eq:time:tai2utc, eq:time:tt, eq:time:ttcent, and
% eq:time:tdb are echoed verbatim in comments in cpp/src/time.cpp and
% cpp/include/star/time.hpp (FR-29 equation-label-to-code traceability);
% renaming them breaks that trace.
\chapter{Time Systems}
\label{ch:time}

\section{Purpose}
\label{sec:time:purpose}

This chapter documents the time-system module (FR-2, decision D-6): the
internal two-part TAI epoch representation, the conversion between UTC
calendar fields and TAI through a bundled, versioned leap-second table, the
exact offset to Terrestrial Time (TT), and the truncated periodic series for
Barycentric Dynamical Time (TDB). Every downstream Phase~2 model consumes
these scales: the Earth-orientation frame chain is a function of UT1 and TT,
and the ephemeris evaluator is a function of TDB. The module is implemented
in \texttt{cpp/src/time.cpp} with its interface in
\texttt{cpp/include/star/time.hpp}; per decision D-2 the core accepts only
numeric calendar fields --- ISO-8601 text parsing lives in the Python
frontend.

\section{Assumptions and domain of validity}
\label{sec:time:domain}

Assumptions:
\begin{enumerate}
  \item \textbf{Constant-offset UTC era only.} UTC is related to TAI by an
    integer number of seconds, which holds from 1972-01-01 onward. The
    pre-1972 ``rubber second'' era (frequency offsets and fractional jumps)
    is out of domain.
  \item \textbf{Leap-second knowledge horizon.} The bundled table encodes
    the published TAI$-$UTC step history through the step effective
    2017-01-01 ($\Delta AT = \qty{37}{\second}$) and is verified against
    IERS Bulletin~C~71 (January 2026), which announced that no leap second
    is introduced at the end of June 2026~\cite{iersbc71}. The first UTC
    date at which a leap second \emph{not} present in the table could take
    effect is therefore 2027-01-01 (the end-of-December-2026 insertion
    opportunity, which Bulletin~C~71 does not rule out).
  \item \textbf{Positive leap seconds only.} Every leap second since 1972
    has been an insertion; the table holds only $+1$ steps and the
    conversion logic renders only inserted seconds (second $60$). A
    negative leap second would require a table and code revision.
  \item \textbf{Proleptic Gregorian calendar.} Calendar dates map to day
    numbers by the integer Julian Day Number algorithm of Fliegel and Van
    Flandern~\cite{fliegel1968}, exact over the module domain (verified
    exhaustively over 1972--2068 at golden generation).
  \item \textbf{Geocentric TDB.} The TDB$-$TT series carries only the
    time-dependent terms; topocentric terms (observer longitude and
    geocentric distance, $<\qty{2.1}{\micro\second}$ at the geoid) are
    omitted, consistent with the D-6 accuracy budget.
\end{enumerate}

Domain of validity: UTC epochs from 1972-01-01 onward; full leap-second
accuracy through 2026-12-31; TDB$-$TT within the D-6 budget of
$\qty{30}{\micro\second}$ over 1600--2200 (section~\ref{sec:time:derivation}).

Out-of-domain behavior, matching the code exactly:
\begin{itemize}
  \item Calendar input before 1972-01-01, an invalid field (month, day,
    hour, minute out of range; non-finite seconds), or second
    ${}\geq 60$ outside an inserted leap second throws
    \texttt{std::domain\_error} (surfaced in Python as
    \texttt{ValueError}). An invalid instant is never silently normalized,
    because normalizing would shift the epoch by up to one second.
  \item Calendar input on or after the table expiry (2027-01-01) is
    \emph{accepted} with the last tabulated offset
    ($\Delta AT = \qty{37}{\second}$). The core exposes the table version
    and expiry programmatically and the Python frontend emits a warning for
    such epochs; the decision is placed in Python because the core never
    reads the clock (D-2) and cannot know whether its table is stale.
\end{itemize}

\section{Derivation}
\label{sec:time:derivation}

\subsection{Time scales}

International Atomic Time (TAI) is the continuous atomic scale, one SI
second per second on the rotating geoid; it has no leap seconds, so every
TAI day spans exactly \qty{86400}{\second}. Coordinated Universal Time
(UTC) differs from TAI by an integer step count,
\begin{equation}
  \mathrm{TAI} = \mathrm{UTC} + \Delta AT(t),
  \label{eq:time:utc2tai}
\end{equation}
where $\Delta AT$ is the tabulated leap-second offset, constant between
steps that take effect at UTC midnights announced by IERS
Bulletin~C~\cite{iersbc71}. Terrestrial Time is defined with a constant
offset from TAI (IAU 1991 Resolution A4, as stated by
Kaplan~\cite{kaplan2005}):
\begin{equation}
  \mathrm{TT} = \mathrm{TAI} + \qty{32.184}{\second}
  \quad\text{exactly.}
  \label{eq:time:tt}
\end{equation}
The J2000 reference epoch is the instant
\begin{equation}
  \text{J2000}
  = \text{2000-01-01T12:00:00.0 TT}
  = \mathrm{JD}\,2451545.0\ \mathrm{TT}
  = \text{2000-01-01T11:59:27.816 TAI},
  \label{eq:time:j2000}
\end{equation}
the TAI clock face following from equation~\eqref{eq:time:tt}.

\subsection{Two-part TAI epoch}

An epoch is stored as the pair $(d, s)$ --- whole TAI days $d$ and TAI
seconds of day $s \in [0, 86400)$ --- counted from 2000-01-01T00:00:00.0
TAI, so J2000 is the pair $(0,\, 43167.816)$ by
equation~\eqref{eq:time:j2000}. Because TT and TAI differ by a constant,
elapsed TT equals elapsed TAI, and the TT Julian-century argument used by
the TDB series and the frame models is
\begin{equation}
  T \;=\; \frac{86400\,d + \bigl(s - 43167.816\bigr)}{86400 \cdot 36525}
  \qquad \text{[Julian centuries since J2000 TT]}.
  \label{eq:time:ttcent}
\end{equation}

The two-part split is the load-bearing design decision of D-6. A single
binary64 count of seconds since J2000 reaches $\sim\!1.9\times10^{9}$ by
2060, where its quantum (one unit in the last place) is
$2^{-52}\times 2^{31} \approx \qty{2.4e-7}{\second}$ --- incompatible with
the $\sim\!\qty{1e-8}{\second}$ of UT1 precision the Earth-rotation-angle
computation requires ($\qty{1e-8}{\second}$ of UT1 is
$\sim\!\qty{7e-12}{\radian}$ of rotation). In the two-part form the day
count is exact and the sub-day part never exceeds \qty{86400}{\second},
whose quantum is $2^{-52} \times 2^{17} = \qty{1.455e-11}{\second}$,
uniform over the whole 2020--2060 span --- three orders of magnitude below
the requirement. Downstream consumers preserve this by differencing or
reducing the parts separately (never summing them first), exactly as the
validation tests compare two-part Julian Dates.

\subsection{Calendar mapping}

Gregorian calendar dates map to integer day numbers by the machine
algorithm of Fliegel and Van Flandern~\cite{fliegel1968}, which computes
the Julian Day Number with truncating integer division,
\begin{equation}
  \begin{aligned}
    a &= \left(m - 14\right)/12,\\
    \mathrm{JDN}(y, m, d) &= d - 32075
      + 1461\,(y + 4800 + a)/4
      + 367\,(m - 2 - 12a)/12\\
      &\quad - 3\,\bigl((y + 4900 + a)/100\bigr)/4,
  \end{aligned}
  \label{eq:time:fvf}
\end{equation}
with every division truncated toward zero; the epoch day count is
$N = \mathrm{JDN} - 2451545$ since $\mathrm{JDN}(\text{2000-01-01}) =
2451545$. The companion inverse algorithm from the same reference recovers
$(y, m, d)$ from $N$. Both directions were verified exhaustively against an
independent implementation for every day from 1972-01-01 through 2068-12-31
at golden generation.

\subsection{UTC to TAI and back, including leap seconds}

For UTC calendar fields $(y, m, d, h, \mathit{min}, s_{\mathrm{utc}})$ the
UTC seconds of day are $\sigma = 3600h + 60\,\mathit{min} +
s_{\mathrm{utc}}$, and by equation~\eqref{eq:time:utc2tai} the epoch is
\begin{equation}
  A = 86400\,N(y,m,d) + \sigma + \Delta AT(y,m,d),
  \qquad
  d_{\mathrm{TAI}} = \lfloor A / 86400 \rfloor,\quad
  s_{\mathrm{TAI}} = A - 86400\,d_{\mathrm{TAI}},
  \label{eq:time:tai2utc}
\end{equation}
where $A$ is TAI seconds since 2000-01-01T00:00:00 TAI. During an inserted
leap second the seconds field satisfies $s_{\mathrm{utc}} \in [60, 61)$, so
$\sigma \in [86400, 86401)$ and the mapping continues smoothly across the
insertion --- no special case is needed in the forward direction beyond
admitting the extended seconds range on exactly the final minute of a day
the table steps after.

For the inverse, observe that the UTC civil day with epoch day number $D$
begins at $A = 86400\,D + \Delta AT(D)$ and ends where the next day begins,
so its TAI duration is $86400 + [\Delta AT(D{+}1) - \Delta AT(D)]$ seconds:
\qty{86401}{\second} for a day ending in an inserted leap second. Given
$(d_{\mathrm{TAI}}, s_{\mathrm{TAI}})$, the containing UTC day is
$d_{\mathrm{TAI}}$ itself when $s_{\mathrm{TAI}} \geq
\Delta AT(d_{\mathrm{TAI}})$ and the preceding day otherwise, and the UTC
seconds of day are $\sigma = A - 86400\,D - \Delta AT(D)$. The value
$\sigma \in [86400, 86401)$ arises exactly when the preceding day ends in
an inserted second and the instant lies inside it; it is rendered as
23:59:60$+$fraction. The round trip is therefore the identity for every
valid UTC instant, including instants inside a leap second --- the property
gated by \testid{time_leap_second_roundtrip}.

\subsection{TDB}

Barycentric Dynamical Time differs from TT by periodic relativistic terms.
The module implements the seven-term truncation of the Fairhead and
Bretagnon series~\cite{fairhead1990} as given by Kaplan~\cite[eq.~2.6]
{kaplan2005}:
\begin{equation}
  \begin{aligned}
    \mathrm{TDB} - \mathrm{TT} \;[\unit{\second}] \;=\;
      &\phantom{+}\,0.001657 \sin(628.3076\,T + 6.2401)\\
      &+ 0.000022 \sin(575.3385\,T + 4.2970)
       + 0.000014 \sin(1256.6152\,T + 6.1969)\\
      &+ 0.000005 \sin(606.9777\,T + 4.0212)
       + 0.000005 \sin(52.9691\,T + 0.4444)\\
      &+ 0.000002 \sin(21.3299\,T + 5.5431)
       + 0.000010\,T \sin(628.3076\,T + 4.2490),
  \end{aligned}
  \label{eq:time:tdb}
\end{equation}
with $T$ from equation~\eqref{eq:time:ttcent} (using TT rather than TDB in
the argument perturbs the result below $\qty{1e-12}{\second}$ and is the
form Kaplan gives). Kaplan quotes $\sim\!\qty{10}{\micro\second}$ accuracy
against the full model over 1600--2200~\cite{kaplan2005}, inside the D-6
budget of \qty{30}{\micro\second}; the observed maximum against the full
Fairhead--Bretagnon model (ERFA \texttt{eraDtdb}) over the golden epochs
spanning 2007--2060 is \qty{5.54e-6}{\second}
(\texttt{tests/golden/time/manifest.toml}). The dominant term is the annual
aberration of the terrestrial clock ($k \sin g$ with $g$ the Earth's mean
anomaly); truncation is acceptable because the ephemeris consumer tolerates
tens of microseconds of argument error (interplanetary velocities of
$\sim\!\qty{40}{\kilo\metre\per\second}$ make \qty{30}{\micro\second}
about \qty{1.2}{\milli\metre} of position).

\section{Discretization and implementation}
\label{sec:time:implementation}

\begin{enumerate}
  \item \textbf{Module and traceability.} The conversions live in
    \texttt{cpp/src/time.cpp}; the source comments echo the labels
    \texttt{eq:time:j2000}, \texttt{eq:time:fvf}, \texttt{eq:time:utc2tai},
    \texttt{eq:time:tai2utc}, \texttt{eq:time:tt}, \texttt{eq:time:ttcent},
    and \texttt{eq:time:tdb} verbatim.
  \item \textbf{Exact integer bookkeeping.} Equations~\eqref{eq:time:fvf}
    and~\eqref{eq:time:tai2utc} are evaluated in 64-bit integer arithmetic
    for every whole-second quantity; the fractional second is split off
    exactly (the fractional part of a binary64 is itself representable) and
    rides through unchanged. The single rounding step in a UTC-to-TAI
    conversion is the final addition
    $s_{\mathrm{TAI}} = \mathrm{fl}(\mathrm{isod} + \mathrm{frac})$, whose
    error is at most $2^{-37} \approx \qty{7.3e-12}{\second}$ (one half ulp
    at \qty{86400}{\second}). Consequently the UTC
    $\rightarrow$ TAI $\rightarrow$ UTC round trip is bit-exact whenever
    the fractional second is representable at that quantum (all golden
    inputs use dyadic fractions), and within
    $\qty{1.5e-11}{\second}$ otherwise.
  \item \textbf{Leap-second table.} The 28-entry TAI$-$UTC step history is
    a \texttt{constexpr} array cited at its definition (IERS Bulletin~C
    series, verified through Bulletin~C~71~\cite{iersbc71}), scanned with a
    fixed-order loop (branch-stable, D-10). The table version string and
    expiry date (2027-01-01) are exposed programmatically
    (\texttt{leap\_table\_info}) so the Python frontend can warn on
    post-expiry epochs; the warning lives in Python because the core never
    reads the clock (D-2).
  \item \textbf{Epoch arithmetic.} \texttt{tai\_add\_seconds} restores the
    $0 \leq s < 86400$ invariant by whole-day reduction plus at most one
    fix-up per side (the floating floor at the boundary can be off by one
    ulp). Precision is $\mathrm{ulp}(\max(|\Delta t|, 86400))$: exact
    nanosecond resolution for within-mission offsets, degrading gracefully
    for offsets of years.
  \item \textbf{Determinism.} All paths are straight-line integer and
    IEEE-754 double arithmetic with fixed evaluation order, compiled with
    fast-math disabled and contraction off (D-10). The golden tests
    exploit this: the committed reference values were generated by a
    Python implementation performing the identical operation sequence,
    and the C++ results are required to match them \emph{bit for bit}
    (libm-dependent quantities --- the TDB sines --- carry a documented
    $\qty{1e-13}{\second}$ tolerance instead).
  \item \textbf{Binding surface.} Epochs cross the pybind11 boundary as
    plain \texttt{(day, sec)} tuples; \texttt{std::domain\_error} maps to
    Python \texttt{ValueError}. The Python mission validator warns (never
    errors) on epochs past the table expiry, so a stale table can never
    block a run, only annotate it.
\end{enumerate}

\section{Validation evidence}
\label{sec:time:validation}

Tables~\ref{tab:time:validation-core} and~\ref{tab:time:validation-py}
list the tests that constitute this chapter's validation evidence. Each
test identifier is machine-checked to exist in the test tree by
\texttt{scripts/lint\_chapters.py}; golden provenance and tolerances are
recorded in \texttt{tests/golden/time/manifest.toml}.

\begin{table}[htbp]
  \centering
  \caption{Validation evidence for the time-system module: core (doctest,
    \texttt{cpp/tests/}).}
  \label{tab:time:validation-core}
  \begin{tabular}{lp{8.2cm}}
    \toprule
    Test ID & Acceptance basis \\
    \midrule
    \testid{time_utc_tai_tt_golden} &
      UTC$\to$TAI$\to$TT at 18 epochs (2007--2060): bit equality with the
      golden reference; two-part agreement with ERFA
      (\texttt{dtf2d}/\texttt{utctai}/\texttt{taitt}) to
      $\leq \qty{1e-9}{\second}$ (Phase~2 exit criterion 1); published SOFA
      cookbook anchor~\cite{sofacookbook}; bit-exact calendar round
      trip. \\
    \testid{time_tdb_series_golden} &
      Equation~\eqref{eq:time:tdb} vs.\ the bit-comparable reference to
      $\leq \qty{1e-13}{\second}$ and vs.\ the full Fairhead--Bretagnon
      model (ERFA \texttt{eraDtdb}) to $\leq \qty{30}{\micro\second}$
      (D-6 budget). \\
    \testid{time_leap_table_golden} &
      All 28 table steps and the day before each vs.\ the published
      history; expiry metadata (2027-01-01); post-expiry lookups return
      \qty{37}{\second}. \\
    \testid{time_leap_second_roundtrip} &
      Bit-exact UTC round trip before, inside, and after the inserted
      second at every boundary since 1972; exactly \qty{2}{\second} of TAI
      elapse from 23:59:59 to 00:00:00 across a leap. \\
    \testid{time_two_part_precision} &
      A \qty{1}{\nano\second} step at epoch 2060 survives the two-part
      representation (and demonstrably vanishes in a single double);
      normalization invariant $0 \leq s < 86400$ under epoch
      arithmetic. \\
    \testid{time_domain_errors} &
      Pre-1972 dates, invalid calendar fields, non-finite seconds, and
      second~60 outside an inserted leap second all throw
      \texttt{std::domain\_error}; Gregorian century rule honored. \\
    \bottomrule
  \end{tabular}
\end{table}

\begin{table}[htbp]
  \centering
  \caption{Validation evidence for the time-system module: binding and
    frontend surface (pytest, \texttt{tests/python/}).}
  \label{tab:time:validation-py}
  \begin{tabular}{lp{6.6cm}}
    \toprule
    Test ID & Acceptance basis \\
    \midrule
    \testid{test_time_bindings_match_goldens} &
      The pybind11 surface reproduces the same goldens (bit equality;
      $\leq \qty{1e-9}{\second}$ vs.\ ERFA). \\
    \testid{test_time_tdb_bindings_golden} &
      TDB series and two-part TDB Julian Date through the bindings. \\
    \testid{test_time_leap_roundtrip_bindings} &
      Leap-second round trips and \texttt{ValueError} domain mapping
      through the bindings; table metadata exposure. \\
    \testid{test_time_two_part_arithmetic_precision} &
      Two-part arithmetic resolves \qty{1}{\nano\second} at 2060 through
      the bindings. \\
    \testid{test_time_epoch_expiry_warning} &
      Mission epochs on or after 2027-01-01 (by UTC calendar date,
      including offset time zones) warn and still validate; earlier epochs
      do not warn. \\
    \bottomrule
  \end{tabular}
\end{table}

\section{References}
\label{sec:time:references}

This chapter relies on Kaplan, USNO Circular No.~179~\cite{kaplan2005} for
the time-scale relationships, the J2000 definition, and the truncated TDB
series; on Fairhead and Bretagnon~\cite{fairhead1990} for the underlying
analytical TDB$-$TT model; on the IERS Bulletin~C series~\cite{iersbc71}
for the leap-second history and its verification horizon; on Fliegel and
Van Flandern~\cite{fliegel1968} for the calendar algorithms; and on the IAU
SOFA cookbook~\cite{sofacookbook} for the published worked example anchoring
the golden vectors. Vallado~\cite{vallado2013} gives the survey treatment
of the same material. Full bibliographic details appear in the
bibliography.
